% !TEX root = main.tex
%=============================================================================
% appendix_theory.tex  --  credit identities and proofs for the shared
% shared transition-table model results (theory content scoped to the kept sections).
%=============================================================================

\section{Proofs for the Shared Transition-Table Model}
\label{app:operator-proofs}

This appendix proves \cref{thm:forward-backward,cor:near-mixing,prop:outcome-nonidentifiability,thm:complete-mixing,prop:loss-path}. All
gradients are Euclidean gradients in transition-probability coordinates. If a
gate occurs more than once, its shared-matrix gradient is the sum of the
occurrence-level gradients.

\subsection{Tangent-space projections}
\label{app:tangent-projections}

Let
\[
 \mathcal T=\{X\in\mathbb R^{K\times K}:X\mathbf1=0\},\qquad
 \mathcal C=\{X\in\mathcal T:\mathbf1^\top X=0\}.
\]
For any matrix \(\bm{G}\),
\begin{equation}
 \operatorname{Proj}_{\mathcal T}(\bm{G})=\bm{G}\bm{\Pi},\qquad
 \operatorname{Proj}_{\mathcal C}(\bm{G})=\bm{\Pi} \bm{G}\bm{\Pi}.
 \label{eq:tangent-projections}
\end{equation}
To verify the first identity, note that \(\bm{G}\bm{\Pi}\) has zero row sums. The
discarded matrix \(GU\) is constant within each row and is orthogonal to every
\(X\in\mathcal T\). Applying the same argument to columns proves the second
identity. This also gives the orthogonal decomposition
\[
 \mathcal T=\mathcal M\oplus\mathcal C,\qquad
 \mathcal M=\{\mathbf1c^\top:\mathbf1^\top c=0\}.
\]
The subspace \(\mathcal M\) changes destination marginals equally for every
source. The subspace \(\mathcal C\) changes source-dependent destinations.
Since a permutation matrix is doubly stochastic, \(\bm{T}_g-\bm{U}\in\mathcal C\).

This is the operator the main text calls \(\Rule{\cdot}\):
\begin{equation}
 \Rule{\bm{G}}\;\define\;\operatorname{Proj}_{\mathcal C}(\bm{G})=\bm{\Pi} \bm{G}\bm{\Pi} .
 \label{eq:rule-operator}
\end{equation}
The centered signals of \cref{eq:centered-signals} are the same projection
applied to vectors. For any probability vector \(\bm{q}\) and any vector \(\bm{b}\),
\[
 \bm{\Pi} \bm{q}=\bm{q}-\bm{u}=\fwdsig,
 \qquad
 \bm{\Pi} \bm{b}=\bm{b}-\operatorname{mean}(\bm{b})\mathbf1=\bwdsig ,
\]
so \(\bm{\Pi}\) removes exactly the part of a signal that is common to every state.
Below we write \(\operatorname{Proj}_{\mathcal C}\) and \(\bm{\Pi}\) directly.

\subsection{Factorization and the depth bound}
\label{app:factorization-proof}

\begin{proof}
\emph{Step 1: differentiate the terminal probability.}
Fix an occurrence of gate $g$ at position $t$ and write $\bm P_t=\bm P_g$ for
its table. By associativity and \cref{eq:forward-backward},
\[
 p_y=\bm{q}_{t-1}^{\top}\bm{P}_t\bm{b}_t=\sum_{i=1}^K\sum_{j=1}^K(\bm q_{t-1})_i(\bm P_t)_{ij}(\bm b_t)_j.
\]
Neither $\bm q_{t-1}$ nor $\bm b_t$ depends on the entries of this particular
occurrence of $\bm P_t$: $\bm q_{t-1}$ is built only from tables at earlier
positions and $\bm b_t$ only from tables at later positions. Differentiating
the display entrywise therefore treats $(\bm q_{t-1})_i$ and $(\bm b_t)_j$ as
constants, giving
\[
 \frac{\partial p_y}{\partial(\bm{P}_t)_{ij}}
 =(\bm{q}_{t-1})_i(\bm{b}_t)_j
 \quad\text{for every }i,j,
 \qquad\text{i.e.}\qquad
 \nabla_{\bm P_t}p_y=\bm q_{t-1}\bm b_t^\top.
\]
Since $\ell_{\rm out}=-\log p_y$, the chain rule gives
$\nabla_{\bm P_t}\ell_{\rm out}=-p_y^{-1}\nabla_{\bm P_t}p_y$, so
\[
 -\nabla_{\bm{P}_t}\ell_{\rm out}
 =\frac{\bm{q}_{t-1}\bm{b}_t^\top}{p_y}.
\]
By \cref{eq:tangent-projections}, projecting a matrix $\bm G$ onto the
feasible tangent space $\mathcal T$ means right-multiplying by $\bm\Pi$; doing
so here, and using that $\bm\Pi$ passes through the scalar $1/p_y$ and acts
only on the second factor of the outer product, $(\bm q_{t-1}\bm b_t^\top)\bm\Pi=\bm q_{t-1}(\bm b_t^\top\bm\Pi)=\bm q_{t-1}(\bm\Pi\bm b_t)^\top$
(using $\bm\Pi^\top=\bm\Pi$), gives
\begin{equation}
 -\operatorname{Proj}_{\mathcal T}\nabla_{\bm{P}_t}\ell_{\rm out}
 =\frac{\bm{q}_{t-1}\bwdsig_t^{\top}}{p_y}.
 \label{eq:feasible-credit}
\end{equation}
Projecting once more onto $\mathcal C$ means left-multiplying \cref{eq:feasible-credit}
by $\bm\Pi$ as well: $\bm\Pi\bigl(\bm q_{t-1}\bwdsig_t^\top\bigr)=(\bm\Pi\bm q_{t-1})\bwdsig_t^\top=\fwdsig_{t-1}\bwdsig_t^\top$,
which gives exactly \cref{eq:conditional-credit}.

\emph{Step 2: characterize vanishing conditional credit.}
A rank-one matrix $\bm x\bm y^\top\in\R^{K\times K}$ is the zero matrix if and
only if $\bm x=\bm0$ or $\bm y=\bm0$: if instead $x_i\ne0$ and $y_j\ne0$ for
some $i,j$, then $(\bm x\bm y^\top)_{ij}=x_iy_j\ne0$. Applying this to
\cref{eq:conditional-credit}, the conditional credit at occurrence $t$
vanishes exactly when $\fwdsig_{t-1}=\bm0$ or $\bwdsig_t=\bm0$. Since
$\bm q_{t-1}$ is a probability vector, $\fwdsig_{t-1}=\bm\Pi\bm q_{t-1}=\bm q_{t-1}-\bm u=\bm0$
exactly when $\bm q_{t-1}=\bm u$, i.e.\ the belief state reaching step $t$ is
exactly uniform. Likewise $\bwdsig_t=\bm b_t-\operatorname{mean}(\bm b_t)\mathbf1=\bm0$
exactly when every entry of $\bm b_t$ agrees with the mean, i.e.\ $\bm b_t$ is
a scalar multiple of $\mathbf1$.

\emph{Step 3: compute the local process gradient and its lower bound.}
For the local target $\ell_t=-\log(\bm P_t)_{r,d}$ with $r=s_{t-1}$, $d=s_t$,
this loss depends on $\bm P_t$ only through the single entry
$(\bm P_t)_{r,d}$, so the identical differentiation as in Step 1 (with
$\bm q_{t-1}\bm b_t^\top$ replaced by the rank-one matrix $\bm e_r\bm e_d^\top$,
which has a $1$ at $(r,d)$ and $0$ elsewhere) gives
$-\nabla_{\bm P_t}\ell_t=\bm e_r\bm e_d^\top/(\bm P_t)_{r,d}$, and hence, by the
same two-sided $\bm\Pi$-multiplication as in Step 1,
\[
 \operatorname{Proj}_{\mathcal C}(-\nabla_{\bm P_t}\ell_t)
 =\frac{(\bm{\Pi} \bm{e}_r)(\bm{\Pi} \bm{e}_d)^\top}{(\bm{P}_t)_{r,d}}.
\]
To bound its Frobenius norm, first note that for any $\bm x,\bm y\in\R^K$,
\[
 \|\bm x\bm y^\top\|_F^2=\sum_{i,j}(x_iy_j)^2
 =\Bigl(\sum_i x_i^2\Bigr)\Bigl(\sum_j y_j^2\Bigr)=\|\bm x\|_2^2\|\bm y\|_2^2,
 \qquad\text{so}\qquad
 \|\bm x\bm y^\top\|_F=\|\bm x\|_2\|\bm y\|_2 .
\]
Next, for any basis vector $\bm e_i$, $\bm\Pi\bm e_i=\bm e_i-\bm u$ has $k$-th
entry $\delta_{ik}-1/K$, so
\[
 \|\bm\Pi\bm e_i\|_2^2
 =\Bigl(1-\tfrac1K\Bigr)^2+(K-1)\Bigl(\tfrac1K\Bigr)^2
 =1-\tfrac2K+\tfrac1{K^2}+\tfrac{K-1}{K^2}
 =1-\tfrac2K+\tfrac1K=1-\tfrac1K,
\]
independently of $i$. Combining the last two displays with $r,d$ in place of
$i$,
\[
 \left\|\operatorname{Proj}_{\mathcal C}(-\nabla_{\bm P_t}\ell_t)\right\|_F
 =\frac{\|\bm\Pi\bm e_r\|_2\|\bm\Pi\bm e_d\|_2}{(\bm P_t)_{r,d}}
 =\frac{1-1/K}{(\bm P_t)_{r,d}}\ge1-\frac1K,
\]
the last inequality because $(\bm P_t)_{r,d}\le1$. This proves
\cref{eq:process-credit-lower}.

\emph{Step 4: the forward and backward residuals contract geometrically.}
Now assume every participating table is doubly stochastic and write
$\bm{P}_j=\bm{U}+\bm{E}_j$ with $\|\bm E_j\|_2\le\rulestr$. Doubly stochastic means
$\bm P_j\mathbf1=\mathbf1$ and $\mathbf1^\top\bm P_j=\mathbf1^\top$; since also
$\bm U\mathbf1=\mathbf1$ and $\mathbf1^\top\bm U=\mathbf1^\top$ (as $\bm U=\mathbf1\mathbf1^\top/K$),
subtracting gives $\bm E_j\mathbf1=\bm0$ and $\mathbf1^\top\bm E_j=\bm0^\top$ for
every participating $j$. Consequently
\[
 \bm U\bm E_j=\tfrac1K\mathbf1(\mathbf1^\top\bm E_j)=\bm0,
 \qquad
 \bm E_j\bm U=\tfrac1K(\bm E_j\mathbf1)\mathbf1^\top=\bm0,
\]
and since $\bm\Pi=\bm I-\bm U$, also $\bm\Pi\bm E_j=\bm E_j-\bm U\bm E_j=\bm E_j$
and $\bm E_j\bm\Pi=\bm E_j-\bm E_j\bm U=\bm E_j$. Transposing both identities
and using $\bm\Pi^\top=\bm\Pi$ shows the same absorption holds for the
transpose: $\bm\Pi\bm E_j^\top=\bm E_j^\top\bm\Pi=\bm E_j^\top$. In
particular $\bm\Pi\bm P_j\bm\Pi=\bm\Pi(\bm U+\bm E_j)\bm\Pi=\bm\Pi\bm U\bm\Pi+\bm\Pi\bm E_j\bm\Pi
=\bm0+\bm E_j=\bm E_j$ (using $\bm\Pi\bm U=\bm0$ and $\bm\Pi\bm E_j\bm\Pi=(\bm\Pi\bm E_j)\bm\Pi=\bm E_j\bm\Pi=\bm E_j$),
so $\bm E_j=\bar{\bm P}_j$ exactly, matching \cref{cor:near-mixing}'s notation.

We claim $\bm{P}_1\cdots \bm{P}_r=\bm{U}+\bm{E}_1\cdots \bm{E}_r$ for every $r\ge1$
(with the empty product at $r=0$ equal to $\bm I$), by induction on $r$. The
base case $r=1$ is the definition $\bm P_1=\bm U+\bm E_1$. For the inductive
step, suppose $\bm P_1\cdots\bm P_{r-1}=\bm U+\bm F$ with
$\bm F=\bm E_1\cdots\bm E_{r-1}$; then $\bm F\mathbf1=\bm E_1\cdots\bm E_{r-2}(\bm E_{r-1}\mathbf1)=\bm0$
by repeated use of $\bm E_j\mathbf1=\bm0$, so $\bm F\bm U=\tfrac1K\bm F\mathbf1\mathbf1^\top=\bm0$.
Then, using $\bm U\bm U=\bm U$ (since $\bm U^2=\tfrac1{K^2}\mathbf1(\mathbf1^\top\mathbf1)\mathbf1^\top=\tfrac1{K^2}\mathbf1\cdot K\cdot\mathbf1^\top=\bm U$),
\[
 \bm P_1\cdots\bm P_r=(\bm U+\bm F)(\bm U+\bm E_r)
 =\bm U\bm U+\bm U\bm E_r+\bm F\bm U+\bm F\bm E_r
 =\bm U+\bm0+\bm0+\bm E_1\cdots\bm E_r,
\]
proving the claim.

Fix an occurrence $t$. Writing $\bm F_{t-1}=\bm E_1\cdots\bm E_{t-1}$ for the
$r=t-1$ instance of the claim, and using $\bm q_{t-1}=(\bm P_1\cdots\bm P_{t-1})^\top\bm e_{s_0}$,
\[
 \bm q_{t-1}=(\bm U+\bm F_{t-1})^\top\bm e_{s_0}=\bm U\bm e_{s_0}+\bm F_{t-1}^\top\bm e_{s_0}=\bm u+\bm F_{t-1}^\top\bm e_{s_0},
\]
using $\bm U^\top=\bm U$ and $\bm U\bm e_{s_0}=\bm u$. Since $\bm\Pi\bm u=\bm u-\bm u=\bm0$,
\[
 \fwdsig_{t-1}=\bm\Pi\bm q_{t-1}=\bm\Pi\bm F_{t-1}^\top\bm e_{s_0}.
\]
Because $\bm F_{t-1}^\top=\bm E_{t-1}^\top\cdots\bm E_1^\top$ and each factor
absorbs a left $\bm\Pi$ (shown above), $\bm\Pi\bm F_{t-1}^\top=\bm F_{t-1}^\top$
exactly, so $\fwdsig_{t-1}=\bm F_{t-1}^\top\bm e_{s_0}$. Writing
$\bm e_{s_0}=\bm u+\bm\Pi\bm e_{s_0}$ and noting
$\bm E_1^\top\bm u=\bm0$ (transpose of $\mathbf1^\top\bm E_1=\bm0^\top$, giving
$\bm E_1^\top\mathbf1=\bm0$, hence $\bm E_1^\top\bm u=\bm E_1^\top\mathbf1/K=\bm0$),
we get $\bm F_{t-1}^\top\bm u=\bm E_{t-1}^\top\cdots\bm E_2^\top(\bm E_1^\top\bm u)=\bm0$,
so
\[
 \fwdsig_{t-1}=\bm F_{t-1}^\top\bm e_{s_0}=\bm F_{t-1}^\top\bm\Pi\bm e_{s_0}.
\]
Taking norms and applying submultiplicativity of the spectral norm across
the $t-1$ factors of $\bm F_{t-1}^\top$,
\begin{equation}
 \|\fwdsig_{t-1}\|_2\le\|\bm E_1^\top\|_2\cdots\|\bm E_{t-1}^\top\|_2\,\|\bm\Pi\bm e_{s_0}\|_2
 =\Bigl(\textstyle\prod_{j<t}\|\bm E_j\|_2\Bigr)\sqrt{1-\tfrac1K},
 \label{eq:forward-tail-bound}
\end{equation}
using $\|\bm E_j^\top\|_2=\|\bm E_j\|_2$ and the norm
$\|\bm\Pi\bm e_{s_0}\|_2=\sqrt{1-1/K}$ from Step 3. The identical argument
applied to $\bm b_t=(\bm E_{t+1}\cdots\bm E_D)\bm e_y+\bm u$ (the same claim
with $r=D-t$ factors $\bm E_{t+1},\ldots,\bm E_D$, applied on the left instead
of transposed, since $\bm b_t=\bm P_{t+1}\cdots\bm P_D\bm e_y$ needs no
transpose) gives
\begin{equation}
 \|\bwdsig_t\|_2\le\Bigl(\textstyle\prod_{j>t}\|\bm E_j\|_2\Bigr)\sqrt{1-\tfrac1K}.
 \label{eq:backward-tail-bound}
\end{equation}
If a common bound $\|\bm E_j\|_2\le\rulestr$ holds for every participating
$j$, \cref{eq:forward-tail-bound,eq:backward-tail-bound} specialize to
$\|\fwdsig_{t-1}\|_2\le\rulestr^{t-1}\sqrt{1-1/K}$ and
$\|\bwdsig_t\|_2\le\rulestr^{D-t}\sqrt{1-1/K}$.

\emph{Step 5: combine.} By \cref{eq:conditional-credit} and
$\|\bm x\bm y^\top\|_F=\|\bm x\|_2\|\bm y\|_2$ again, and using
$\bm E_j=\bar{\bm P}_j$ from Step 4 and $p_y\ge\alpha$,
\[
 \left\|\operatorname{Proj}_{\mathcal C}(-\nabla_{\bm{P}_t}\ell_{\rm out})\right\|_F
 =\frac{\|\fwdsig_{t-1}\|_2\|\bwdsig_t\|_2}{p_y}
 \le\frac{1-1/K}{\alpha}\prod_{j\ne t}\|\bar{\bm P}_{g_j}\|_2
 =\frac{1-1/K}{\alpha}\,\chi_t,
\]
which is \cref{eq:general-credit-bound}. If additionally
$\|\bar{\bm P}_j\|_2\le\rulestr$ for every participating $j$, then
$\chi_t\le\rulestr^{(t-1)+(D-t)}=\rulestr^{D-1}$ regardless of $t$, giving
\cref{eq:depth-credit-bound}. This proves \cref{cor:near-mixing}.
\end{proof}

\paragraph{A comparison in one normalization.} \Cref{eq:process-credit-lower}
and \cref{eq:depth-credit-bound} are both per-occurrence bounds, so contrasting
$1-1/K$ against $\rulestr^{D-1}/\alpha$ already compares like with like; the
population objectives are not automatically comparable, since $L_{\rm proc}$
in \cref{eq:shared-process-loss} averages $D$ such occurrences with an
explicit $1/D$ while $L_{\rm out}$ does not. Carrying the per-occurrence
process bound through that average does not, however, introduce a hidden
$D^{-1}$ decay: \cref{thm:complete-mixing} computes the resulting population
credit exactly as $f_g(\bm{Q}_g-\bm{U})$, with $f_g=D^{-1}\sum_t\Pr(g_t=g)$ the
gate's own visiting \emph{frequency}, not a count --- $\Theta(1)$ in $D$
whenever a gate's per-step probability does not itself vanish with $D$, as
in this paper's tasks. The shared-table outcome gradient instead sums
occurrences without a $1/D$ factor:
\[
 \nabla_{\bm P_g}L_{\rm out}
 =\mathbb E\!\left[\sum_{t=1}^D\mathbf1\{g_t=g\}
      \nabla_{\bm P_t}\ell_{\rm out}\right].
\]
Linearity of $\Rule{\cdot}$, the triangle inequality, and
\cref{cor:near-mixing}, with the same $\alpha$ and $\rulestr$ for every
example, therefore give
\begin{equation}
 \left\|\Rule{-\nabla_{\bm P_g}L_{\rm out}}\right\|_F
 \le D f_g\,\frac{1-1/K}{\alpha}\,\rulestr^{D-1}.
 \label{eq:population-outcome-credit-bound}
\end{equation}
For clean supervision at complete mixing, process credit is
$f_g(\bm T_g-\bm U)$; for fixed $K$ and $f_g$ bounded below independently
of $D$, its norm is $\Theta(1)$. With fixed $\alpha>0$, the population
outcome ceiling is $O(D f_g\rulestr^{D-1})$, still exponentially suppressed
for fixed $\rulestr<1$ despite the polynomial prefactor $D$. The
per-occurrence ceiling remains $O(\rulestr^{D-1})$, and the exact zero at
complete mixing is unchanged. These are upper bounds on outcome credit,
not matching lower bounds, since particular products can cancel further.

Outcome supervision can fail to identify the true rule for a second,
independent reason beyond the vanishing-gradient geometry just proved: a
source--gate pair that is never visited on any supported trajectory leaves
no trace in $L_{\rm out}$ at all, regardless of gradient magnitude.

\begin{proposition}[Outcome non-identifiability under under-excitation]
\label{prop:outcome-nonidentifiability}
Let $\bm{P}_g=\bm{T}_g$ for every gate except possibly at one row $i^\star$ of
one gate $g^\star$, where $\bm{P}_{g^\star}(i^\star,\cdot)=\bm{v}$ for an
arbitrary row-stochastic $\bm{v}$. If the pair $(g^\star,i^\star)$ is never
visited on the true trajectory of any prompt in the training distribution's
support --- i.e., $(g_t,s_{t-1})\ne(g^\star,i^\star)$ for every $t\in[D]$ and
every $(s_0,g_{1:D})$ with positive probability --- then $L_{\rm out}$ is
unchanged by this replacement for every such $\bm{v}$: the modified table is
a global minimizer of $L_{\rm out}$ whenever $\bm{P}_g=\bm{T}_g$ everywhere
is, so $L_{\rm out}$ alone cannot identify $\bm{T}_{g^\star}(i^\star,\cdot)$
from any amount of training data drawn from this distribution.
\end{proposition}
\begin{proof}
Since every other row of every $\bm{P}_g$ already equals $\bm{T}_g$'s and each
$\bm{T}_g$ is a permutation, induction on $t$ shows the model's own belief
state $\bm{q}_t=\bm{P}_{g_t}^\top\cdots\bm{P}_{g_1}^\top\bm{e}_{s_0}$ equals
$\bm{e}_{s_t}$ exactly, for the true trajectory $s_t$, at every step of
every supported prompt: $\bm{q}_0=\bm{e}_{s_0}$ by definition, and if
$\bm{q}_{t-1}=\bm{e}_{s_{t-1}}$ with $(g_t,s_{t-1})\ne(g^\star,i^\star)$ by
hypothesis, then $\bm{q}_t=\bm{P}_{g_t}^\top\bm{e}_{s_{t-1}}=\bm{T}_{g_t}^\top
\bm{e}_{s_{t-1}}=\bm{e}_{s_t}$ regardless of $\bm{v}$, since row $i^\star$ of
$\bm{P}_{g^\star}$ is never the row selected. Row $i^\star$ is therefore
never read by the forward computation of any supported prompt, so
$p_y=\bm{e}_{s_0}^\top\bm{P}_{g_1}\cdots\bm{P}_{g_D}\bm{e}_y$ --- and hence
$L_{\rm out}$, an expectation over this same distribution --- does not
depend on $\bm{v}$.
\end{proof}

\Cref{prop:outcome-nonidentifiability}'s ambiguity closes under full
coverage: if every $(g,i)$ pair is visited with positive probability, the
argument no longer applies. A second, independent source of
non-identifiability survives even then.

\begin{theorem}[Group-symmetry non-identifiability]
\label{thm:group-symmetry-nonidentifiability}
Let $\{\bm{T}_g\}_{g\in\mathcal G}\subset S_K$ be the true gate permutations
and let $\bm A\in S_K$, $\bm A\ne\bm I$, commute with every $\bm T_g$ (an
element of the centralizer of the generated group
$\langle\bm T_g:g\in\mathcal G\rangle$), with $\bm A^D=\bm I$ for the task's
fixed depth $D$. Define $\bm T'_g:=\bm T_g\bm A$ for every $g\in\mathcal G$.
Then $\bm T'_g\ne\bm T_g$ for every $g$, yet
\[
 \bm T'_{g_1}\cdots\bm T'_{g_D}=\bm T_{g_1}\cdots\bm T_{g_D}
\]
for every sequence $g_1,\ldots,g_D\in\mathcal G$ (repeats allowed), so
$\{\bm T'_g\}$ and $\{\bm T_g\}$ are distinct gate-to-permutation
assignments inducing identical terminal behavior on every depth-$D$ word.
Consequently $L_{\rm out}$ does not identify $\{\bm T_g\}$ from
$\{\bm T'_g\}$ under any training distribution supported on depth-$D$
prompts, however complete its coverage, whereas process supervision (which
displays $s_t=\bm T_{g_t}(s_{t-1})$ under the fixed labeling that defines
$\{\bm T_g\}$) distinguishes them at every step.
\end{theorem}
\begin{proof}
$\bm T'_g=\bm T_g\iff\bm T_g\bm A=\bm T_g\iff\bm A=\bm I$ (multiply by
$\bm T_g^{-1}$), contradicting $\bm A\ne\bm I$; so $\bm T'_g\ne\bm T_g$ for
every $g$. For the composition, since $\bm A$ commutes with every $\bm T_g$,
\[
 \bm T'_{g_1}\cdots\bm T'_{g_D}
 =(\bm T_{g_1}\bm A)(\bm T_{g_2}\bm A)\cdots(\bm T_{g_D}\bm A)
 =\bm T_{g_1}\bm T_{g_2}\cdots\bm T_{g_D}\,\bm A^D
 =\bm T_{g_1}\cdots\bm T_{g_D},
\]
moving each $\bm A$ rightward past every subsequent $\bm T_{g_j}$ one
factor at a time (each such transposition uses
$\bm A\bm T_{g_j}=\bm T_{g_j}\bm A$) and using $\bm A^D=\bm I$. Both
assignments therefore give the same $p_y=\bm e_{s_0}^\top\bm T_{g_1}\cdots
\bm T_{g_D}\bm e_y$ for every prompt, hence the same $L_{\rm out}$ for every
distribution over depth-$D$ prompts, so no amount of outcome-only training
data distinguishes them. A process target instead displays $s_t$ directly
under the $\{\bm T_g\}$-labeling at every step, so a model matching
$\{\bm T'_g\}$ at row $i$ is directly penalized whenever
$\bm T'_g(i)=\bm T_g(\bm A(i))\ne\bm T_g(i)$, which holds exactly when
$\bm A(i)\ne i$ (since $\bm T_g$ is injective) --- i.e.\ at every row $\bm A$
moves, and every row unless $\bm A$ happens to fix it.
\end{proof}

Whether such an $\bm A$ exists depends entirely on the gate family: a
generating set rich enough to produce all or most of $S_K$ has trivial
centralizer (the centralizer of $S_n$ in itself is trivial for $n\ge3$), so
this specific construction gives nothing for a near-universal family like
this paper's Boolean-circuit gates. It is not vacuous, however: any abelian
generated group has this property with room to spare, since every element
commutes with every other. \Cref{app:shift-symmetry-task} constructs one
such family explicitly, as a worked example rather than a trained-network
measurement.

A trivial centralizer of $\langle\bm T_g\rangle$ itself is not the right
sufficient condition for identifiability, however: a gate family can have
trivial centralizer and still admit a distinct zero-loss executor at a
fixed depth $D$, since the escaping symmetry need only commute with the
\emph{differences} between gates, not with each gate individually. A
sufficient criterion uses the difference group.

\Cref{thm:fixed-depth-identifiability} is stated in \cref{sec:operator-credit}.
This appendix gives the proof.
\begin{proof}
\emph{Step 1: every $\bm P_g$ is a permutation.} Full support gives
$\bm P_{g_1}\cdots\bm P_{g_D}=\bm T_{g_1}\cdots\bm T_{g_D}=:\bm T$, a
permutation, for every word. Taking determinants, every $\bm P_{g_i}$ is
invertible ($\det\bm T=\pm1\ne0$ forces every factor's determinant nonzero).
Isolating one factor, $\bm P_{g_1}^{-1}=(\bm P_{g_2}\cdots\bm P_{g_D})\bm T^{-1}$
is a product of nonnegative matrices, hence nonnegative; since $\bm P_{g_1}$
is also nonnegative (row-stochastic), it is a nonnegative matrix with a
nonnegative inverse, hence a monomial matrix, and row-stochasticity forces
its single nonzero entry per row to equal $1$: $\bm P_{g_1}$ is a genuine
permutation matrix. The same argument isolates every other factor of every
word, so every $\bm P_g:=\bm S_g$ used anywhere is a permutation.

\emph{Step 2: two one-sided decompositions of $\bm S_g$.} Fix a reference
gate $h_0$. Full support gives, for every $g\in\mathcal G$,
\[
 \bm S_g\bm S_{h_0}^{D-1}=\bm T_g\bm T_{h_0}^{D-1}
 \quad\text{and}\quad
 \bm S_{h_0}^{D-1}\bm S_g=\bm T_{h_0}^{D-1}\bm T_g,
\]
the first from words $(g,h_0,\ldots,h_0)$, the second from words
$(h_0,\ldots,h_0,g)$ (each with $D-1$ copies of $h_0$). Since every
$\bm S_g,\bm T_g$ is invertible, these give
\[
 \bm S_g=\bm T_g\bm A,\qquad \bm A:=\bm T_{h_0}^{D-1}\bm S_{h_0}^{-(D-1)},
 \qquad\qquad
 \bm S_g=\bm B\bm T_g,\qquad \bm B:=\bm S_{h_0}^{-(D-1)}\bm T_{h_0}^{D-1},
\]
for every $g\in\mathcal G$, with $\bm A,\bm B$ independent of $g$.

\emph{Step 3: $\bm B\in C_{S_K}(H)$.} Equating the two decompositions,
$\bm T_g\bm A=\bm B\bm T_g$, i.e.\ $\bm A=\bm T_g^{-1}\bm B\bm T_g$, for
every $g$. Since the left side does not depend on $g$, for any
$g_1,g_2\in\mathcal G$, $\bm T_{g_1}^{-1}\bm B\bm T_{g_1}=\bm T_{g_2}^{-1}
\bm B\bm T_{g_2}$, and rearranging,
$\bm B(\bm T_{g_2}\bm T_{g_1}^{-1})=(\bm T_{g_2}\bm T_{g_1}^{-1})\bm B$: $\bm B$
commutes with every generator of $H$, so $\bm B\in C_{S_K}(H)$.

\emph{Step 4: conclude.} If $C_{S_K}(H)=\{\bm I\}$, then $\bm B=\bm I$, so
$\bm T_g=\bm S_g\bm B^{-1}\cdot\bm B=\bm B\bm T_g$ gives (from Step 2's second
relation) $\bm S_g=\bm B\bm T_g=\bm T_g$ directly, for every $g\in\mathcal G$.
\end{proof}

This gives a complementary sufficient identifiability criterion to
\cref{thm:group-symmetry-nonidentifiability}'s ambiguity construction.
Any alternative zero-loss executor induces a nonidentity
$\bm B\in C_{S_K}(H)$; the proof does not show that every such centralizer
element constructs an alternative at the chosen depth.
Since $H\le\langle\bm T_g\rangle$ and this inclusion can be strict,
$C_{S_K}(\langle\bm T_g\rangle)\le C_{S_K}(H)$; triviality of the former
does not imply triviality of the latter in general. The sufficient
criterion therefore concerns differences between gates.
\Cref{app:identifiability-certificate} verifies $C_{S_K}(H)=\{\bm I\}$
computationally for this paper's own Boolean-circuit gate family, so its
outcome failure is not attributable to an alternative exact factorization
of the terminal map.

\subsection{Population stationarity}
\label{app:population-stationarity}
\begin{proof}
Throughout, fix the gate sequence $g_{1:D}$ and reason about the randomness
in $s_0$ conditional on it; by hypothesis $s_0$ is uniform and independent
of $g_{1:D}$.

\emph{Step 1: the true trajectory is uniform at every step.} Each $\bm T_g$
is a permutation matrix, so for the fixed sequence $g_{1:D}$ the map
$s_0\mapsto s_t$ obtained by applying $\bm T_{g_1},\ldots,\bm T_{g_t}$ in
turn is a composition of $t$ permutations of $[K]$, hence itself a
permutation of $[K]$. A permutation applied to a uniform random variable is
again uniform, so conditional on $g_{1:D}$, each $s_t$ ($t=0,\ldots,D$) is
uniform on $[K]$.

\emph{Step 2: the outcome credit vanishes identically for $t<D$.} Set
$\bm P_g=\bm U$ for every gate. Fix an occurrence at position $t<D$; by
\cref{eq:forward-backward}, $\bm b_t=\bm P_{t+1}\cdots\bm P_D\bm e_y=\bm U\cdots\bm U\bm e_y$,
a product of at least one $\bm U$ applied (possibly repeatedly) to $\bm e_y$.
Since $\bm U\bm x=\mathbf1(\mathbf1^\top\bm x)/K$ is always a multiple of
$\mathbf1$ for any vector $\bm x$, and $\bm U\bm u=\bm u$ (because
$\mathbf1^\top\bm u=1$), applying $\bm U$ once already gives $\bm u$, and
applying it again leaves $\bm u$ unchanged; so $\bm b_t=\bm u$ regardless of
how many factors of $\bm U$ remain in the product. Hence
$\bwdsig_t=\bm\Pi\bm b_t=\bm\Pi\bm u=\bm0$, and by \cref{eq:conditional-credit}
the projected credit at this occurrence,
$\Rule{-\nabla_{\bm P_t}\ell_{\rm out}}=\fwdsig_{t-1}\bwdsig_t^\top/p_y$, is
the zero matrix --- for every realization of the prompt, not merely in
expectation.

\emph{Step 3: the outcome credit at $t=D$ vanishes only after averaging over
$s_D$.} At the final occurrence, $\bm q_{D-1}=\bm P_{D-1}^\top\cdots\bm P_1^\top\bm e_{s_0}=\bm U^\top\cdots\bm U^\top\bm e_{s_0}=\bm u$
by the same collapsing argument (using $\bm U^\top=\bm U$), and
$\bm b_D=\bm e_y=\bm e_{s_D}$ (the empty product). Then
$p_y=\bm q_{D-1}^\top\bm P_D\bm b_D=\bm u^\top\bm U\bm e_{s_D}=\bm u^\top\bm u=1/K$,
using $\bm U\bm e_{s_D}=\bm u$ and $\mathbf1^\top\bm u=1$ so
$\bm u^\top\bm u=(1/K)\mathbf1^\top\bm u=1/K$. By \cref{eq:feasible-credit},
the tangent-space (not yet fully $\mathcal C$-projected) credit at this
occurrence is
\[
 -\operatorname{Proj}_{\mathcal T}\nabla_{\bm P_D}\ell_{\rm out}
 =\frac{\bm q_{D-1}\bwdsig_D^\top}{p_y}
 =\frac{\bm u(\bm\Pi\bm e_{s_D})^\top}{1/K}
 =K\bm u(\bm\Pi\bm e_{s_D})^\top
 =\mathbf1(\bm\Pi\bm e_{s_D})^\top,
\]
using $K\bm u=\mathbf1$. For a \emph{fixed} realization of $s_D$ this need
not be zero: the outcome objective can retain a nonzero tangent-space
signal at any single instance. But taking the expectation over $s_D$
conditional on $g_{1:D}$, and using Step 1 ($s_D$ uniform given $g_{1:D}$,
so $\E[\bm e_{s_D}\mid g_{1:D}]=\bm u$),
\[
 \E\bigl[\bm\Pi\bm e_{s_D}\mid g_{1:D}\bigr]=\bm\Pi\,\E[\bm e_{s_D}\mid g_{1:D}]=\bm\Pi\bm u=\bm0,
\]
so the $t=D$ contribution to the \emph{population} feasible gradient
vanishes as well. Combining Steps 2 and 3, and summing over any repeated
occurrences of the same gate across positions (\cref{app:operator-proofs}'s
convention that a shared table's gradient is the sum of its occurrence-level
gradients), gives $(-\nabla_{\bm P_g}L_{\rm out})\bm\Pi=\bm0$ for every gate
$g$: the first identity in \cref{eq:process-recovers-table}.

\emph{Step 4: the process credit equals $f_g(\bm T_g-\bm U)$ exactly.} Fix
an occurrence of gate $g$ at position $t$, and let $r=s_{t-1}$ be the
(random) source state reached there. At $\bm P_t=\bm U$ every entry equals
$1/K$, including the true local target $(\bm P_t)_{r,\bm T_g(r)}=1/K$, so by
the computation in \cref{app:factorization-proof} the projected local-process
credit at this occurrence is
\[
 \operatorname{Proj}_{\mathcal C}(-\nabla_{\bm P_t}\ell_t)
 =\frac{(\bm{\Pi} \bm{e}_r)(\bm{\Pi} \bm{e}_{\bm{T}_g(r)})^\top}{1/K}
 =K(\bm{\Pi} \bm{e}_r)(\bm{\Pi} \bm{e}_{\bm{T}_g(r)})^\top.
\]
By Step 1, conditional on $g_{1:D}$ the source $r$ is uniform on $[K]$, so
averaging over $r$,
\begin{align*}
 \frac1K\sum_{r=1}^K K(\bm{\Pi} \bm{e}_r)(\bm{\Pi} \bm{e}_{\bm{T}_g(r)})^\top
 &=\bm{\Pi}\left(\sum_{r=1}^K\bm e_r\bm e_{\bm{T}_g(r)}^\top\right)\bm{\Pi}\\
 &=\bm{\Pi} \bm{T}_g\bm{\Pi},
\end{align*}
where the first equality distributes $\bm\Pi$ out of the finite sum by
linearity, and the second uses that $\sum_r\bm e_r\bm e_{\bm T_g(r)}^\top$ is
exactly the entrywise definition of the permutation matrix $\bm T_g$ (its
$(r,j)$ entry is $1$ iff $j=\bm T_g(r)$). Since $\bm T_g$ is doubly
stochastic, the same computation as for $\bm E_j$ in
\cref{app:factorization-proof} gives $\bm U\bm T_g=\bm T_g\bm U=\bm U$
(concretely, $\bm U\bm T_g=\tfrac1K\mathbf1(\mathbf1^\top\bm T_g)=\tfrac1K\mathbf1\mathbf1^\top=\bm U$,
using $\mathbf1^\top\bm T_g=\mathbf1^\top$), so
\[
 \bm\Pi\bm T_g\bm\Pi=(\bm I-\bm U)\bm T_g(\bm I-\bm U)
 =\bm T_g-\bm U\bm T_g-\bm T_g\bm U+\bm U\bm T_g\bm U
 =\bm T_g-\bm U-\bm U+\bm U=\bm T_g-\bm U.
\]
Averaging further over the $D$ positions of the sequence, and summing
repeated occurrences of gate $g$ as before, the definition
$f_g=D^{-1}\sum_{t=1}^D\Pr(g_t=g)$ gives exactly
\[
 \Rule{-\nabla_{\bm P_g}L_{\rm proc}}=f_g(\bm T_g-\bm U),
\]
the second identity in \cref{eq:process-recovers-table}.

\emph{Step 5: the row maximizer recovers $\bm T_g$.} Row $i$ of $\bm T_g$ is
the one-hot row vector with a $1$ at column $\bm T_g(i)$, so row $i$ of
$\bm T_g-\bm U$ has value $1-1/K$ at column $\bm T_g(i)$ and value $-1/K$ at
every other column. Since $1-1/K>-1/K$ for every $K\ge1$, the unique row
maximizer of $f_g(\bm T_g-\bm U)$'s row $i$, whenever $f_g>0$, is column
$\bm T_g(i)$; this holds for every $i$, so the row-wise maximizer of the
process rule credit identifies every successor of $\bm T_g$.
\end{proof}

\subsection{Loss on the affine table path}
\label{app:affine-path-proof}

\begin{proposition}[Order of improvement along the true rules]
\label{prop:loss-path}
Let each $\bm{T}_g$ be a permutation matrix and set $\bm{P}_g(a)=\bm{U}+a(\bm{T}_g-\bm{U})$ for
$0\le a\le1$. For any prompt distribution and depth $D\ge2$,
\begin{align}
 L_{\rm out}(a)&=\log K-\log\!\left(1+(K-1)a^D\right),
 \label{eq:outcome-flat-path}\\
 L_{\rm proc}(a)&=\log K-\log\!\left(1+(K-1)a\right).
 \label{eq:process-linear-path}
\end{align}
At $a=0$, the first $D-1$ outcome derivatives vanish, while
$L_{\rm proc}'(0)=-(K-1)$.
\end{proposition}
\begin{proof}
Fix any gates $g_1,\ldots,g_D$ and write $\bm E_g=\bm T_g-\bm U$, so
$\bm P_g(a)=\bm U+a\bm E_g$ by definition.

\emph{Step 1: $\bm U\bm E_g=\bm E_g\bm U=\bm0$, and
$\bm E_g\bm E_h=\bm T_g\bm T_h-\bm U$ for any two gates $g,h$.} Since $\bm T_g$
is a permutation matrix it is doubly stochastic:
$\bm T_g\mathbf1=\mathbf1$, $\mathbf1^\top\bm T_g=\mathbf1^\top$; and
$\bm U\mathbf1=\mathbf1$, $\mathbf1^\top\bm U=\mathbf1^\top$ as well. Subtracting
gives $\bm E_g\mathbf1=\bm0$ and $\mathbf1^\top\bm E_g=\bm0^\top$, exactly as
in \cref{app:factorization-proof}, from which
$\bm U\bm E_g=\tfrac1K\mathbf1(\mathbf1^\top\bm E_g)=\bm0$ and
$\bm E_g\bm U=\tfrac1K(\bm E_g\mathbf1)\mathbf1^\top=\bm0$. For the product
identity, expand directly:
\[
 \bm E_g\bm E_h=(\bm T_g-\bm U)(\bm T_h-\bm U)
 =\bm T_g\bm T_h-\bm T_g\bm U-\bm U\bm T_h+\bm U\bm U.
\]
Since $\bm T_g$ is doubly stochastic, the same computation shows
$\bm T_g\bm U=\tfrac1K(\bm T_g\mathbf1)\mathbf1^\top=\tfrac1K\mathbf1\mathbf1^\top=\bm U$,
and symmetrically $\bm U\bm T_h=\bm U$; and $\bm U\bm U=\bm U$ since
$\bm U^2=\tfrac1{K^2}\mathbf1(\mathbf1^\top\mathbf1)\mathbf1^\top=\tfrac1{K^2}\mathbf1\cdot K\cdot\mathbf1^\top=\bm U$.
Substituting,
\[
 \bm E_g\bm E_h=\bm T_g\bm T_h-\bm U-\bm U+\bm U=\bm T_g\bm T_h-\bm U.
\]

\emph{Step 2: the product formula, by induction on $D$.} We claim
\[
 \prod_{t=1}^D\bm P_{g_t}(a)=\prod_{t=1}^D(\bm U+a\bm E_{g_t})
 =\bm U+a^D\bigl(\bm T_{g_1}\cdots\bm T_{g_D}-\bm U\bigr)
\]
for every $D\ge1$. The base case $D=1$ is the definition
$\bm P_{g_1}(a)=\bm U+a\bm E_{g_1}=\bm U+a^1(\bm T_{g_1}-\bm U)$. For the
inductive step, suppose the claim holds at $D-1$, i.e.\
$\prod_{t=1}^{D-1}\bm P_{g_t}(a)=\bm U+a^{D-1}\bm F$ with
$\bm F:=\bm T_{g_1}\cdots\bm T_{g_{D-1}}-\bm U$. Since a product of
permutation matrices is again a permutation matrix, Step 1 applies to $\bm F$
exactly as it did to each $\bm E_g$, giving $\bm U\bm F=\bm F\bm U=\bm0$.
Then, using $\bm U\bm U=\bm U$,
\begin{align*}
 \prod_{t=1}^D\bm P_{g_t}(a)
 &=\bigl(\bm U+a^{D-1}\bm F\bigr)\bigl(\bm U+a\bm E_{g_D}\bigr)\\
 &=\bm U\bm U+a\,\bm U\bm E_{g_D}+a^{D-1}\bm F\bm U+a^D\bm F\bm E_{g_D}\\
 &=\bm U+\bm0+\bm0+a^D\bm F\bm E_{g_D},
\end{align*}
using $\bm U\bm E_{g_D}=\bm0$ (Step 1 applied to $g_D$). Finally, the same
product identity as in Step 1 --- which only used that both factors are
built from doubly stochastic matrices minus $\bm U$ --- gives
$\bm F\bm E_{g_D}=(\bm T_{g_1}\cdots\bm T_{g_{D-1}}-\bm U)(\bm T_{g_D}-\bm U)=\bm T_{g_1}\cdots\bm T_{g_D}-\bm U$,
proving the claim at $D$.

\emph{Step 3: reading off the terminal and local entries.} Fix a prompt
$(s_0,g_{1:D})$ with true answer $y$, so that $\bm T_{g_1}\cdots\bm T_{g_D}$'s
$(s_0,y)$ entry equals $1$ (this composed permutation matrix's row $s_0$ has
its single $1$ exactly at column $y$, the true terminal state). By
\cref{eq:terminal-probability} and Step 2,
\[
 p_y(a)=\bm e_{s_0}^\top\Bigl[\bm U+a^D(\bm T_{g_1}\cdots\bm T_{g_D}-\bm U)\Bigr]\bm e_y
 =\frac1K+a^D\Bigl(1-\frac1K\Bigr)=\frac{1+(K-1)a^D}{K},
\]
using $\bm e_{s_0}^\top\bm U\bm e_y=1/K$ (every entry of $\bm U$ equals $1/K$)
and $\bm e_{s_0}^\top(\bm T_{g_1}\cdots\bm T_{g_D})\bm e_y=1$ (the $(s_0,y)$
entry just identified). The right-hand side does not depend on which prompt
was fixed, only on $a$, $K$, $D$; since $L_{\rm out}(a)=\E[-\log p_y(a)]$ and
$-\log p_y(a)=\log K-\log(1+(K-1)a^D)$ identically for every prompt in the
support, the expectation over prompts equals this same constant:
\[
 L_{\rm out}(a)=\log K-\log\bigl(1+(K-1)a^D\bigr),
\]
proving \cref{eq:outcome-flat-path}.

For the process loss, fix any single occurrence of a gate $g$ with true
local transition $r\to\bm T_g(r)$. By definition $\bm P_g(a)=\bm U+a(\bm T_g-\bm U)$,
and row $r$ of $\bm T_g$ places a $1$ at column $\bm T_g(r)$, so
\[
 (\bm P_g(a))_{r,\bm T_g(r)}=\frac1K+a\Bigl(1-\frac1K\Bigr)=\frac{1+(K-1)a}{K},
\]
again independent of $r$, $g$, and the prompt. Every one of the $D$ local
terms in $L_{\rm proc}(a)=D^{-1}\sum_t\E[-\log(\bm P_{g_t}(a))_{s_{t-1},s_t}]$
therefore equals the same constant $\log K-\log(1+(K-1)a)$, and so does
their average:
\[
 L_{\rm proc}(a)=\log K-\log\bigl(1+(K-1)a\bigr),
\]
proving \cref{eq:process-linear-path}.

\emph{Step 4: the derivative claims.} Using the Taylor expansion
$\log(1+z)=z-\tfrac12z^2+O(z^3)$ as $z\to0$: setting $z=(K-1)a^D$ in the
outcome-loss identity,
\[
 L_{\rm out}(a)=\log K-(K-1)a^D+O(a^{2D}),
\]
whose derivatives of order $1,\ldots,D-1$ at $a=0$ all vanish, since every
nonconstant term is $O(a^D)$ or higher order; the $D$-th derivative at $a=0$
is $-(K-1)D!$, nonzero. Setting $z=(K-1)a$ in the process-loss identity,
\[
 L_{\rm proc}(a)=\log K-(K-1)a+O(a^2),\qquad\text{so}\qquad L_{\rm proc}'(0)=-(K-1).
\]
Since the identities in Step 3 hold for every individual prompt (not merely
after averaging), and each population loss is obtained by averaging the
\emph{same} function of $a$ over the prompt distribution, both
\cref{eq:outcome-flat-path,eq:process-linear-path} and the derivative
claims hold regardless of how the initial state $s_0$ is distributed.
\end{proof}

\subsection{Reliability mixtures and gradient contributions}
\label{app:mixture-proof}

The constructions above concern deterministic correct targets.  To model the
reliability experiment, instead let $X$ be a prompt, let $B$ be an independent
Bernoulli variable with success probability $\rho$, and let the trace be drawn
from the clean or corrupt condition according to $B$.  The final answer remains
correct.  Within the corrupt condition, the trace may also be random.

\paragraph{Gradient of a reliability mixture.}
Let $\mathcal L_{\rm clean}(\bm{\theta};X)$ and
$\mathcal L_{\rm corrupt}(\bm{\theta};X,Z)$ be differentiable per-example losses,
where $Z$ describes trace corruption.  Suppose the data laws are independent of
$\bm{\theta}$ and differentiation commutes with expectation, as it does for finite
support at differentiable parameter values.  Use a fixed per-example weighting
rule.  Then
\begin{equation}
 L_\rho=\rho L_{\rm clean}+(1-\rho)L_{\rm corrupt},\qquad
 \nabla L_\rho=\rho g_{\rm clean}+(1-\rho)g_{\rm corrupt},
 \label{eq:mixture-gradient}
\end{equation}
where $g_b=\nabla L_b$ at the same parameter value. This is an elementary
consequence of the law of total expectation, used repeatedly below. We
record it once here rather than re-deriving it at each use.

\begin{proof}
Write the per-example loss as
$\mathcal L(\bm\theta;X,Z,B)=B\,\mathcal L_{\rm clean}(\bm\theta;X)+(1-B)\,\mathcal L_{\rm corrupt}(\bm\theta;X,Z)$,
so that $\mathcal L=\mathcal L_{\rm clean}$ when $B=1$ (probability $\rho$) and
$\mathcal L=\mathcal L_{\rm corrupt}$ when $B=0$ (probability $1-\rho$),
matching the description of the reliability experiment. Taking the
expectation over $B$ first, with $(X,Z)$ held fixed, and then over $(X,Z)$,
\begin{align*}
 L_\rho&=\E_{X,Z,B}[\mathcal L(\bm\theta;X,Z,B)]\\
 &=\E_{X,Z}\Bigl[\E_B[B]\,\mathcal L_{\rm clean}(\bm\theta;X)+\E_B[1-B]\,\mathcal L_{\rm corrupt}(\bm\theta;X,Z)\Bigr]\\
 &=\rho L_{\rm clean}+(1-\rho)L_{\rm corrupt},
\end{align*}
using $\E_B[B]=\rho$, $\E_B[1-B]=1-\rho$, and that $B$ is independent of
$(X,Z)$ so it may be integrated out first. This is the first identity in
\cref{eq:mixture-gradient}. Because $\rho$ and $1-\rho$ do not depend on
$\bm\theta$, and differentiation commutes with the expectations defining
$L_{\rm clean}$ and $L_{\rm corrupt}$ by hypothesis,
\begin{align*}
 \nabla L_\rho&=\nabla\bigl[\rho L_{\rm clean}(\bm\theta)+(1-\rho)L_{\rm corrupt}(\bm\theta)\bigr]\\
 &=\rho\nabla L_{\rm clean}(\bm\theta)+(1-\rho)\nabla L_{\rm corrupt}(\bm\theta)
 =\rho g_{\rm clean}+(1-\rho)g_{\rm corrupt}
\end{align*}
by linearity of the gradient operator, where both $g_{\rm clean}$ and
$g_{\rm corrupt}$ are evaluated at the \emph{same} $\bm\theta$: the identity
is silent about gradients computed at different parameter values, so it
does not license substituting in gradients from two separately trained
models.
\end{proof}

\proofstep{Normalization and finite datasets.}
The proposition uses per-example averaging.  If loss is instead normalized by
the total number of supervised tokens, the effective clean weight depends on
mean token counts: in a population token average it is
$\rho\,\E[T_{\rm clean}]/(\rho\E[T_{\rm clean}]+(1-\rho)\E[T_{\rm corrupt}])$.
The intended weight $\rho$ is recovered when those means agree.  In a fixed
finite dataset the exact identity uses the empirical clean fraction and the
losses on the assigned subsets.  Reusing nested assignments does not make every
subset identical across $\rho$.

\proofstep{From a fixed-checkpoint sign change to an actual threshold.}
For a fixed reference direction $v$ with $a=\ip{g_{\rm clean}}{v}>0>b=\ip{g_{\rm corrupt}}{v}$,
the mixture projection $\ip{\nabla L_\rho}{v}=\rho a+(1-\rho)b$ changes sign at
some $\rho_\star\in(0,1)$. By itself this is only a fixed-checkpoint fact,
silent about training dynamics. \Cref{prop:noisy-process-credit} below
computes $a$ and $b$ explicitly for the corruption law these experiments
use, turning this observation into the actual recovery threshold.

\proofstep{Between-condition gradient variation.}
If a sampled gradient $\bm{G}$ has conditional means $\bm{g}_{\rm clean},\bm{g}_{\rm corrupt}$
and conditional covariance matrices $\bm{\Sigma}_{\rm clean},\bm{\Sigma}_{\rm corrupt}$,
then the law of total covariance gives
\begin{align}
 \operatorname{Cov}(\bm{G})
 &=\rho\bm{\Sigma}_{\rm clean}+(1-\rho)\bm{\Sigma}_{\rm corrupt}\\
 &\quad+\rho(1-\rho)(g_{\rm clean}-g_{\rm corrupt})
                         (g_{\rm clean}-g_{\rm corrupt})^\top.
 \label{eq:mixture-covariance}
\end{align}
To verify the last term, the conditional means differ from the mixture mean by
$(1-\rho)(g_{\rm clean}-g_{\rm corrupt})$ and
$-\rho(g_{\rm clean}-g_{\rm corrupt})$, respectively.  Weighting their outer
products gives the stated coefficient $\rho(1-\rho)$.  Thus reliability can
change both the mean gradient and its variability.

Assigning corruption per whole trace (matching the released generators,
\cref{app:legacy-protocol}) versus per occurrence changes this covariance
without changing the mean gradient, since both share the same
$\rho$-weighted average. We do not carry that comparison further here, as
it concerns gradient variance rather than the reliability threshold of
\cref{cor:noise-threshold}, which depends only on the mean.


\subsection{The reliability threshold for local credit}
\label{app:noise-threshold-proof}

This subsection proves
\cref{prop:noise-convergence,cor:noise-threshold,eq:noisy-process-credit},
in the shared transition-table model of \cref{sec:operator-credit}. The
convergence and recovery-threshold claims hold for any row-stochastic
corruption law $\bm{C}_g$; only the closed-form gradient identity at
$\bm{P}_g=\bm{U}$ additionally needs \cref{thm:complete-mixing}'s
hypotheses (uniform $s_0$ independent of $g_{1:D}$) and double
stochasticity of every $\bm{C}_g$. We
first record the clean ($\rho=1$) special case, which the general argument
below extends unchanged to noisy targets.

\begin{remark}[Clean convergence, the $\rho=1$ case]
\label{cor:clean-convergence}
Let each $\bm{T}_g$ be a permutation matrix and $f_g>0$. In row-logit coordinates
$\bm{P}_{g,i,\cdot}=\softmax(a_{g,i})$, the population process loss $L_{\rm proc}$
is convex in each $a_{g,i}$ and separable across the $KM$ pairs $(g,i)$. A
clean target places all its mass on the true successor, so a sufficiently
small constant step size drives gradient-descent loss values to their
infimum from any finite initialization, with every predicted row
$\bm{P}_{g,i,\cdot}$ converging to the one-hot distribution $\bm{T}_g[i,\cdot]$
(unattained at finite parameters, since $\bm{T}_g$ is a permutation). Once
every row favors the true successor, induction gives exact greedy
execution. No comparable guarantee holds for $L_{\rm out}$:
\cref{thm:complete-mixing} shows its projected gradient is exactly zero at
the natural starting point $\bm{P}_g=\bm{U}$, so at that point there is no
first-order rule-directed signal toward $\bm{T}_g$.
\end{remark}

\proofstep{What the corruption law is.}
The released generators build a corrupted trace from its own displayed prefix:
at step $t$ they compute the valid successor $\Phi(\widetilde s_{t-1},g_t)$ of
the \emph{corrupted} state and then emit a different state, so every displayed
transition is invalid under its own prefix rather than merely different from the
clean trajectory.  Write $\bm{C}_g$ for the resulting conditional law of the
displayed successor given the displayed source and the gate.  A
$\rho$-reliable pool then induces the conditional local law
\begin{equation}
 \bm{Q}_g=\rho \bm{T}_g+(1-\rho)\bm{C}_g ,
 \label{eq:induced-local-law}
\end{equation}
provided the displayed source has the same law in the clean and the corrupted
condition.  It does whenever $\bm{C}_g$ is doubly stochastic: the initial state is
uniform, every $\bm{T}_g$ is a permutation, and a doubly stochastic kernel preserves
the uniform law, so the displayed state is uniform at every step of both walks.
Assigning corruption to whole traces rather than to single steps therefore
changes how displayed pairs are correlated \emph{within} a trace, but not this
conditional law, and only the conditional law enters the population gradient.
For a general, non-doubly-stochastic $\bm{C}_g$, a whole-trace $\rho$-reliable
pool need not induce \cref{eq:induced-local-law} exactly, since the
displayed source's own law can then depend on whether the trace was
corrupted; \cref{def:local-reliability} sidesteps this by taking
\cref{eq:induced-local-law} itself as the definition of the local
reliability $\rho$, which coincides with the whole-trace rate exactly in
the doubly-stochastic case just verified.
The two assignments do change the \emph{variance} of the aggregated gradient,
through the within-condition terms of \cref{app:mixture-proof}'s general
mixture-variance identity (\cref{eq:mixture-covariance}); that appendix does
not carry the per-trace-versus-per-occurrence comparison further, since
neither the mean gradient nor the reliability threshold below depends on it.

\begin{proposition}[Population process credit under corrupted targets]
\label{prop:noisy-process-credit}
In the setting above, with every $\bm{C}_g$ doubly stochastic, at $\bm{P}_g=\bm{U}$
\begin{equation}
 \Rule{-\nabla_{\bm{P}_g}L_{\rm proc}}
 =f_g\,\bm{\Pi} \bm{Q}_g\bm{\Pi}
 =f_g\left(\bm{Q}_g-\bm{U}\right),
 \label{eq:noisy-credit-proof}
\end{equation}
and the row-wise maximizer of $\Rule{-\nabla_{\bm{P}_g}L_{\rm proc}}$ equals the
row-wise maximizer of $\bm{Q}_g$.  If a corrupted target is uniform on the $K-1$
states other than the valid successor, then
$\bm{C}_g-\bm{U}=-\frac{1}{K-1}(\bm{T}_g-\bm{U})$ and
\begin{equation}
 \Rule{-\nabla_{\bm{P}_g}L_{\rm proc}}
 =f_g\,\frac{K\rho-1}{K-1}\left(\bm{T}_g-\bm{U}\right),
 \label{eq:symmetric-credit-proof}
\end{equation}
so the credit is positively aligned with the true rule iff $\rho>1/K$, zero at
$\rho=1/K$, and negatively aligned below.  If instead every corrupted target
follows one fixed rule $\bm{T}_g'$ with $\bm{T}_g'(i)\ne \bm{T}_g(i)$, the row maximizer is
correct iff $\rho>1/2$.
\end{proposition}

\begin{proof}
\emph{Step 1: one displayed pair, at $\bm P_g=\bm U$.} For an occurrence of
gate $g$ with displayed source $i$ and displayed target $j$, the per-example
process term is $\ell=-\log(\bm P_g)_{ij}$. This is exactly the local target
considered in \cref{app:factorization-proof} with $r=i$, $d=j$, so its
feasible conditional gradient is, by that proof,
\[
 \Rule{-\nabla_{\bm P_g}\ell}=\frac{(\bm\Pi\bm e_i)(\bm\Pi\bm e_j)^\top}{(\bm P_g)_{ij}}.
\]
At $\bm P_g=\bm U$, every entry $(\bm P_g)_{ij}=1/K$, so this equals
$K(\bm\Pi\bm e_i)(\bm\Pi\bm e_j)^\top$.

\emph{Step 2: average over the displayed target $j$, then the displayed
source $i$.} Conditional on the displayed source $i$, the displayed target
$j$ has law $\bm Q_g(i,\cdot)$ by \cref{eq:induced-local-law} --- a mixture
of the true successor (probability $\rho$) and the corrupted successor
(probability $1-\rho$); this is \cref{eq:mixture-gradient}'s
mixture-linearity applied to the indicator vector $\bm e_j$ in place of a
whole loss. Since $\E[\bm e_j\mid i]=\bm Q_g(i,\cdot)^\top$ (the row
$\bm Q_g(i,\cdot)$ read as a column vector), averaging Step 1's expression
over $j$ gives $K(\bm\Pi\bm e_i)\bigl(\bm\Pi\,\bm Q_g(i,\cdot)^\top\bigr)^\top$
for each fixed $i$.

The displayed source $i$ is itself uniform on $[K]$ (the same argument as
Step 1 of \cref{app:population-stationarity}, which used only that each
$\bm T_g$ is a permutation and $s_0$ is uniform, both still assumed here),
so averaging further over $i$,
\[
 \frac1K\sum_{i=1}^{K}K\,\bm{\Pi} \bm{e}_i\bigl(\bm\Pi\,\bm Q_g(i,\cdot)^\top\bigr)^\top
 =\bm{\Pi}\Bigl(\sum_{i=1}^K \bm{e}_i\,\bm Q_g(i,\cdot)\Bigr)\bm{\Pi}
 =\bm{\Pi} \bm{Q}_g\bm{\Pi},
\]
where the first equality distributes $\bm\Pi$ out of the finite sum by
linearity, and the second uses that $\sum_i\bm e_i\bm Q_g(i,\cdot)$
reassembles $\bm Q_g$ row by row. Averaging further over the $D$ positions
of the sequence contributes the visitation frequency $f_g$, exactly as in
\cref{eq:process-recovers-table}, giving
$\Rule{-\nabla_{\bm P_g}L_{\rm proc}}=f_g\,\bm{\Pi} \bm{Q}_g\bm{\Pi}$.

It remains to simplify $\bm\Pi\bm Q_g\bm\Pi$. Since
$\bm{Q}_g=\rho \bm{T}_g+(1-\rho)\bm{C}_g$ is a convex combination of two doubly
stochastic matrices, it is itself doubly stochastic:
$\bm Q_g\mathbf1=\rho\bm T_g\mathbf1+(1-\rho)\bm C_g\mathbf1=\rho\mathbf1+(1-\rho)\mathbf1=\mathbf1$,
and symmetrically $\mathbf1^\top\bm Q_g=\mathbf1^\top$. As in
\cref{app:factorization-proof}, this gives
$\bm U\bm Q_g=\tfrac1K\mathbf1(\mathbf1^\top\bm Q_g)=\tfrac1K\mathbf1\mathbf1^\top=\bm U$
and symmetrically $\bm Q_g\bm U=\bm U$, so
\[
 \bm{\Pi} \bm{Q}_g\bm{\Pi}
 =(\bm I-\bm U)\bm Q_g(\bm I-\bm U)
 =\bm{Q}_g-\bm{U}\bm{Q}_g-\bm{Q}_g\bm{U}+\bm{U}\bm{Q}_g\bm{U}
 =\bm{Q}_g-\bm{U}-\bm U+\bm U=\bm{Q}_g-\bm{U},
\]
using $\bm U\bm Q_g\bm U=\bm U(\bm Q_g\bm U)=\bm U\bm U=\bm U$. This proves
\cref{eq:noisy-credit-proof}. Because $\bm U$ adds the same constant $1/K$
to every entry of a row, subtracting it from $\bm Q_g$ shifts every entry of
that row by the same amount and so does not change which column achieves
the row maximum: the row-wise maximizer of
$\Rule{-\nabla_{\bm{P}_g}L_{\rm proc}}=f_g(\bm{Q}_g-\bm{U})$ (whenever $f_g>0$)
coincides with the row-wise maximizer of $\bm Q_g$ itself.

\emph{Step 3: the two named corruption laws.} For the symmetric law, a
corrupted target is uniform on the $K-1$ wrong states, so row $i$ of
$\bm C_g$ places mass $1/(K-1)$ on every $j\ne\bm T_g(i)$ and $0$ on
$\bm T_g(i)$; equivalently
$\bm{C}_g=\frac{1}{K-1}(\mathbf1\mathbf1^\top-\bm{T}_g)=\frac{K\bm U-\bm{T}_g}{K-1}$
(using $\mathbf1\mathbf1^\top=K\bm U$). Then
\[
 \bm{C}_g-\bm{U}
 =\frac{K\bm U-\bm{T}_g}{K-1}-\bm U
 =\frac{K\bm U-\bm{T}_g-(K-1)\bm{U}}{K-1}
 =\frac{\bm U-\bm T_g}{K-1}
 =-\frac{1}{K-1}(\bm{T}_g-\bm{U}).
\]
Substituting into $\bm Q_g-\bm U=\rho(\bm T_g-\bm U)+(1-\rho)(\bm C_g-\bm U)$
(which follows from $\bm Q_g=\rho\bm T_g+(1-\rho)\bm C_g$ by subtracting
$\bm U=\rho\bm U+(1-\rho)\bm U$ from both sides),
\[
 \bm{Q}_g-\bm{U}
 =\rho(\bm{T}_g-\bm{U})-\frac{1-\rho}{K-1}(\bm T_g-\bm U)
 =\Bigl[\rho-\tfrac{1-\rho}{K-1}\Bigr](\bm{T}_g-\bm{U})
 =\frac{K\rho-1}{K-1}(\bm{T}_g-\bm{U}),
\]
where the last step multiplies numerator and denominator by $K-1$ and
simplifies $\rho(K-1)-(1-\rho)=K\rho-1$. Since $\bm T_g\ne\bm U$ for
$K\ge2$ (a permutation matrix is not uniform), this coefficient's sign
determines the alignment of the credit with $\bm T_g-\bm U$: positive for
$\rho>1/K$, zero at $\rho=1/K$, and negative for $\rho<1/K$, proving
\cref{eq:symmetric-credit-proof} and the stated trichotomy.

For a coherent wrong rule $\bm{T}_g'$ disagreeing with $\bm T_g$ everywhere,
row $i$ of $\bm C_g$ places all its mass on the single state
$\bm T_g'(i)\ne\bm T_g(i)$, so row $i$ of
$\bm Q_g=\rho\bm T_g+(1-\rho)\bm C_g$ places mass $\rho$ on $\bm T_g(i)$,
mass $1-\rho$ on $\bm T_g'(i)$, and $0$ elsewhere. The row maximizer is
$\bm{T}_g(i)$ exactly when $\rho>1-\rho$, i.e.\ $\rho>1/2$.
\end{proof}

\proofstep{What gradient descent reaches.}
\Cref{eq:noisy-credit-proof} is a statement at one parameter.  In this
substrate it extends to the solution actually reached, for any
row-stochastic $\bm{C}_g$, not only doubly stochastic ones. Write
$\pi_{g,i}:=D^{-1}\sum_t\Pr(g_t=g,\text{ displayed source}_{t}=i)$ for the
frequency with which gate $g$ is applied at displayed source $i$; under
\cref{thm:complete-mixing}'s hypotheses and doubly stochastic $\bm{C}_g$, the
displayed source is uniform on $[K]$ independent of $g$ (as in
\cref{prop:noisy-process-credit}'s proof), giving the special case
$\pi_{g,i}=f_g/K$ used there, but the argument below needs only
$\pi_{g,i}>0$ at each row it claims recovered. Positive gate frequency
$f_g>0$ alone does not give this: a corruption law that is not doubly
stochastic can make some displayed sources at gate $g$ arbitrarily rare
or entirely unreached, so row coverage $\pi_{g,i}>0$ is the relevant,
possibly strictly stronger, condition in general. In row-logit
coordinates $\bm{P}_{g,i,\cdot}=\softmax(a_{g,i})$, the population process loss is
\[
 L_{\rm proc}
 =-\sum_{g,i}\pi_{g,i}\sum_{j=1}^{K}
   Q_{g,i,j}\log\softmax(a_{g,i})_j ,
\]
a sum of softmax cross-entropies that is convex in each $a_{g,i}$ and separable
across the $KM$ pairs $(g,i)$, with minimizing distribution
$\bm{P}_{g,i,\cdot}=\bm{Q}_{g,i,\cdot}$ whenever $\pi_{g,i}>0$. Each row Hessian is
$\pi_{g,i}(\operatorname{diag}(p)-pp^\top)$ and has operator norm at most
$\pi_{g,i}/2$. A sufficiently small constant step size, for example
$0<\eta\le1/\max_{g,i}\pi_{g,i}$, therefore gives convergence of gradient-descent
loss values to their infimum from any finite initialization. The excess loss
is $\sum_{g,i}\pi_{g,i}\KL(\bm{Q}_{g,i,\cdot}\|\bm{P}_{g,i,\cdot})$, so the predicted
rows with $\pi_{g,i}>0$ approach $\bm{Q}_g(i,\cdot)$. When $\bm{Q}_g$ has zero entries, this limit can require
unbounded logits; a finite logit minimizer is not asserted. If every true
successor is a strict row maximum of $\bm{Q}_g$ at every row with $\pi_{g,i}>0$, convergence eventually makes
all decoded rules correct, and induction gives exact execution. This argument
concerns the tabular population objective; the finite-sample fit minimizes a
different, empirical objective.

\proofstep{A fully corrupted corpus still determines the executor.}
At $\rho=0$ with symmetric corruption, $\bm{C}_g(i,j)=\frac{1}{K-1}
\one\{j\ne \bm{T}_g(i)\}$, so $\bm{T}_g(i)$ is the unique $j$ with $\bm{C}_g(i,j)=0$: the law
of the displayed pairs determines $\bm{T}_g$ exactly.  The cross-entropy optimum is
$\bm{P}_g=\bm{C}_g$, whose row maximizer is never $\bm{T}_g(i)$, and by
\cref{eq:symmetric-credit-proof} the credit is $-\frac{f_g}{K-1}(\bm{T}_g-\bm{U})$, anti-aligned
with the true rule at $\frac{1}{K-1}$ of the clean magnitude.  Identifiability
of the executor from a corpus and recovery of it by this objective are thus
different properties, and the \textsc{Corrupted} condition of
\cref{tab:supervision-comparison} separates them;
\cref{app:noise-threshold} measures both.

\proofstep{Finite samples.}
The population statement requires the true successor to carry the largest
share of a row of $\bm{Q}_g$; the finite-sample statement requires it to win a
plurality of that row's counts.  With $N$ circuits of depth $D$ there are
$ND/(KM)$ displayed pairs per cell on average under uniform gate-string
sampling.

\begin{proof}[Proof of \cref{prop:finite-sample-recovery}]
Fix $(g,i)$, write $T=\bm{T}_g(i)$, $p_j=\bm{Q}_g(i,j)$,
$p_\star=\max_{j\ne T}p_j=p_T-\Delta_{g,i}$, $N=N_{g,i}$, and $\Delta=\Delta_{g,i}$,
and let $n_j$ be the empirical count of successor $j$ among $N$ i.i.d.\ draws
from $\bm{Q}_g(i,\cdot)$. Each $n_j$ is marginally $\mathrm{Binomial}(N,p_j)$,
so Hoeffding's inequality gives $\Pr[n_T\le N(p_T-\Delta/2)]\le\exp(-N\Delta^2/2)$,
and for every $j\ne T$, since $p_j\le p_\star$,
\[
 \Pr\bigl[n_j\ge N(p_\star+\Delta/2)\bigr]
 \ \le\ \Pr\bigl[n_j\ge N(p_j+\Delta/2)\bigr]
 \ \le\ \exp(-N\Delta^2/2).
\]
If $n_T>N(p_T-\Delta/2)$ and $n_j<N(p_\star+\Delta/2)$ for every $j\ne T$, then
since $p_\star+\Delta/2=p_T-\Delta/2$, every $n_j<n_T$ and the empirical
plurality at row $(g,i)$ is $T$. A union bound over the $K-1$ wrong states
and the lower-tail event gives
$\Pr[\text{row }(g,i)\text{ fails}]\le K\exp(-N_{g,i}\Delta_{g,i}^2/2)$; a
further union bound over the $KM$ cells, each budgeted failure probability
$\delta/(KM)$, gives \cref{eq:finite-sample-recovery} and the simultaneous
guarantee.
\end{proof}

Sampling error can prevent recovery even above $1/K$, particularly
when the population margin is small, and \cref{eq:finite-sample-recovery}
is a sufficient condition, not a tight one: it says nothing about which
rows fail when it is violated, or about how per-cell counts $N_{g,i}$
concentrate around their mean $ND/(KM)$ under the sampler actually used.
For fixed $\rho>1/K$ and positive gate
frequencies, empirical row frequencies converge to $\bm{Q}_g$ as the number of
independently sampled circuits grows; finite-sample recovery need not vary
monotonically with sample size or reliability.  We report the measured
finite-sample frontier in \cref{app:noise-threshold} rather than rely on
this bound alone, and the count form of the same statement --- a fixed
number of valid traces against a growing number of corrupted ones, with the
population direction surviving while corrupted traces number strictly fewer
than $(K-1)$ times the valid ones --- is measured there too.

\proofstep{Scope.}
This is a population and finite-sample fact about the tabular objective
alone. The measured gap between $1/K$ and the trained-network frontier of
\cref{fig:legacy-reliability} is a separate quantity, reported rather than
derived here.

\subsection{A one-dimensional diagnostic under a clean terminal answer}
\label{app:outcome-rescue-diagnostic}

This subsection gives the full statement, proof, and figure behind
\cref{sec:outcome-rescue}'s diagnostic. \Cref{cor:noise-threshold} analyzes
the process objective in isolation, but the paper's own process condition
never actually presents a noisy trace alone: every process example
supervises a possibly corrupted trace \emph{together with} the
always-correct final answer (\cref{sec:formats,sec:reliability}). We ask
here, as a one-dimensional diagnostic and nothing more, what that
combination does to the threshold, restricting to symmetric corruption and
the affine path $\bm{P}_g(a)=\bm{U}+a(\bm{T}_g-\bm{U})$ of
\cref{prop:loss-path}. This restriction is not a simplification we could
lift by substituting \cref{cor:noise-threshold}'s general
corruption-structure functional \cref{eq:general-corruption-threshold}: along the affine path every row of
$\bm{P}_g(a)$ places the same mass $(1-a)/K$ on \emph{every} wrong
successor, by construction, so no point on this path can distinguish a
strong competitor from a weak one --- exactly the distinction
$\kappa(\bm{C})=\max_{g,i}(c_\star-c_T)$ needs. A one-dimensional path that
also tracked competing-mass concentration would need an additional
coordinate moving toward a specific wrong successor, not just away from
the uniform table. \Cref{prop:outcome-rescue} and \cref{cor:noise-threshold}
are complementary for this reason, not two special cases of one formula.
Along this path the noisy local-target loss is,
exactly,
\begin{equation}
 L^\rho_{\rm tr}(a) = \log K - \rho\log\bigl(1+(K-1)a\bigr) - (1-\rho)\log(1-a),
 \label{eq:noisy-process-path-loss}
\end{equation}
and \cref{prop:loss-path} gives the clean terminal-answer loss
$L_{\rm out}(a)=\log K-\log(1+(K-1)a^D)$, the loss of a predictor with no
displayed trace to condition on, marginalizing over $D$ compositions from
$s_0$. The paper's own process examples do not present this task at the
answer position: the displayed trace before it is entirely clean or
entirely corrupted (\cref{sec:reliability}). When clean, the token just
before the answer is already the true final state, so predicting the
answer is a \emph{copy} of the immediately preceding displayed token, a
task we take to draw on a copying capability rather than the shared
transition-table competence $a$, and whose loss $L_{\rm copy}$ we
accordingly treat as independent of $a$. When corrupted, the displayed
trace ends at a corrupted state unrelated to the true answer, so recovering
it requires the same $D$-fold composition from the (uncorrected) prompt as
$L_{\rm out}(a)$. The answer position's loss is therefore the mixture
$\rho L_{\rm copy}+(1-\rho)L_{\rm out}(a)$, and since $L_{\rm copy}$ does not
depend on $a$ it contributes only an additive constant, dropped below. For a
relative weight $\beta\ge0$ on the terminal answer, define the combined
objective
\begin{equation}
 J_{\rho,\beta}(a) := L^\rho_{\rm tr}(a) + \beta(1-\rho)\, L_{\rm out}(a).
 \label{eq:combined-loss}
\end{equation}
This is closer to the paper's actual per-example process loss than simply
weighting $L_{\rm out}(a)$ by a $\rho$-independent $\beta$, though it still
rests on treating $L_{\rm copy}$ as $a$-independent. Separately, the
paper's own per-character training loss masks every character of the
target string, not only the $D$ state positions and the final answer: gate
specifiers and separators are supervised too. These structural characters
are copied verbatim from the prompt and never corrupted, so their
per-character loss does not depend on $\rho$ or $a$ either and drops out of
$dJ/da$ the same way. Restricted to the state and answer characters ---
encoded in the same $\log_2K$ bits in this paper's tasks --- the mask
weights the $D$ noisy positions against the $1$ clean position in a $D:1$
ratio, which is what corresponds to $\beta=1/D$. This equal-width
coincidence, not an even split over $D+1$ positions, is what that part of
the correspondence rests on.

\begin{proposition}[Reliability frontier under a clean terminal answer]
\label{prop:outcome-rescue}
Let each $\bm{T}_g$ be a permutation matrix, $D\ge2$, $0\le a<1$, and let
$\beta\ge0$ satisfy, at every $(a,D)$ considered,
\begin{equation}
 K-\beta\, C(a,D)>0,
 \label{eq:rescue-beta-condition}
\end{equation}
with $C(a,D)$ as defined in \cref{eq:rescue-correction} below --- the exact
condition under which the denominator of \cref{eq:rescue-threshold} stays
strictly positive. A simple sufficient condition, proved below and
independent of $a$ and $D$, is
\begin{equation}
 0\le\beta<\frac{2}{K-1};
 \label{eq:rescue-beta-sufficient}
\end{equation}
the paper's own $\beta=1/D$ falls outside this sufficient range at small
$D$ (e.g.\ $1/2>2/15$ at $D=2$, $K=16$), but is checked below, and
numerically in \cref{fig:prop8-frontier}, to satisfy
\cref{eq:rescue-beta-condition} directly at $K=16$ for every $D$ from $2$
to at least $64$, where the denominator stays above $8.48$. This is a
claim about \emph{where the derivative of one fixed diagnostic objective
changes sign along one fixed path}, not about the trajectory of a trained
network under arbitrary SGD: along \cref{eq:combined-loss},
$J_{\rho,\beta}'(a)<0$ --- i.e., moving further along this path lowers the
objective --- exactly when $\rho>\rho_c(a,D,\beta)$, where, writing
\begin{equation}
 C(a,D):=\frac{D(K-1)a^{D-1}(1-a)\bigl(1+(K-1)a\bigr)}{1+(K-1)a^D},
 \label{eq:rescue-correction}
\end{equation}
\begin{equation}
 \rho_c(a,D,\beta):=\frac{1+(K-1)a-\beta\, C(a,D)}{K-\beta\, C(a,D)}.
 \label{eq:rescue-threshold}
\end{equation}
At $a=0$, $C(0,D)=0$, so \cref{eq:rescue-beta-condition} holds trivially
and $\rho_c(0,D,\beta)=1/K$ for every $\beta\ge0$: at complete mixing the
terminal-answer term contributes nothing regardless of its weight,
recovering \cref{cor:noise-threshold}'s threshold exactly. At $\beta=0$,
\cref{eq:rescue-threshold} reduces to the process-only threshold
$\frac{1+(K-1)a}{K}$ used in \cref{cor:noise-threshold}'s proof; for
$\beta>0$ and $a\in(0,1)$, $\rho_c(a,D,\beta)$ is strictly smaller, so the
clean answer never raises the reliability an oracle-directed step needs
\emph{relative to not having it at all}, at the same competence $a$ ---
whether the threshold rises or falls as $a$ itself increases is a separate
question, addressed below.
\end{proposition}
\begin{proof}
$D(K-1)a^{D-1}(1-a)\le\frac{K-1}2$ for every $a\in[0,1]$, $D\ge2$: writing
$f(a):=Da^{D-1}(1-a)$, $f'(a)=0$ at $a^\star=(D-1)/D$, where
$f(a^\star)=(1-1/D)^{D-1}$, a sequence in $D\ge2$ maximized at $D=2$
(value $\tfrac12$) and decreasing thereafter, so $f(a)\le\tfrac12$
always. Since also $\frac{1+(K-1)a}{1+(K-1)a^D}\le K$ (numerator at most
$K$, denominator at least $1$, as $a\le1$), $C(a,D)\le\frac{K(K-1)}2$ for
every $a,D$. Hence $\beta<\frac2{K-1}$ (\cref{eq:rescue-beta-sufficient})
gives $\beta\,C(a,D)<\frac2{K-1}\cdot\frac{K(K-1)}2=K$, i.e.\
\cref{eq:rescue-beta-condition}, so \cref{eq:rescue-beta-sufficient} is
indeed sufficient for the assumed \cref{eq:rescue-beta-condition}, though
not necessary: the paper's own $\beta=1/D$ falls outside the sufficient
range \cref{eq:rescue-beta-sufficient} at small $D$ (e.g.\ $1/2>2/15$ at
$D=2$, $K=16$) yet satisfies \cref{eq:rescue-beta-condition} directly ---
verified numerically above by evaluating $K-\beta C(a,D)$ across
$a\in[0,1)$ for $K=16$ and every $D$ from $2$ to $64$, where it stays
above $8.48$ --- which is all the proposition requires. If instead
$\beta C(a,D)>K$ at some $(a,D)$, i.e.\ \cref{eq:rescue-beta-condition}
fails there, the argument below still identifies the sign change, but
dividing by the now-negative $K-\beta C(a,D)$ reverses the inequality, so
$J_{\rho,\beta}'(a)<0$ would hold for $\rho<\rho_c(a,D,\beta)$ instead; we
do not use or need this regime, since \cref{eq:rescue-beta-condition} is
assumed throughout.

Differentiating \cref{eq:noisy-process-path-loss} as in
\cref{cor:noise-threshold}'s proof gives
$\frac{dL^\rho_{\rm tr}}{da}=\frac{1-K\rho+(K-1)a}{(1-a)[1+(K-1)a]}$, and
\cref{prop:loss-path}'s proof gives
$\frac{dL_{\rm out}}{da}=-\frac{D(K-1)a^{D-1}}{1+(K-1)a^D}$. By
\cref{eq:combined-loss},
\[
 J_{\rho,\beta}'(a)
 =\frac{1-K\rho+(K-1)a}{(1-a)[1+(K-1)a]}
 -\beta(1-\rho)\,\frac{D(K-1)a^{D-1}}{1+(K-1)a^D}.
\]
Both denominators are positive for $0\le a<1$, so, multiplying through by
$(1-a)(1+(K-1)a)(1+(K-1)a^D)>0$, $J_{\rho,\beta}'(a)<0$ iff
\[
 \bigl(1-K\rho+(K-1)a\bigr)\bigl(1+(K-1)a^D\bigr)
 <\beta(1-\rho)\,D(K-1)a^{D-1}(1-a)\bigl(1+(K-1)a\bigr),
\]
i.e., in the notation of \cref{eq:rescue-correction},
$1-K\rho+(K-1)a<\beta(1-\rho)C(a,D)$. Expanding
$\beta(1-\rho)C(a,D)=\beta C(a,D)-\beta C(a,D)\rho$ and collecting the
$\rho$ terms gives
$1+(K-1)a-\beta C(a,D)<\rho\bigl(K-\beta C(a,D)\bigr)$. Under the assumed
\cref{eq:rescue-beta-condition}, $K-\beta C(a,D)>0$, so dividing by it
gives exactly \cref{eq:rescue-threshold}.
At $a=0$ and $D\ge2$, $a^{D-1}=0$, so $C(0,D)=0$ and
$\rho_c(0,D,\beta)=1/K$. For $\beta=0$ the correction vanishes entirely,
giving the process-only threshold; for $\beta>0$ and $a\in(0,1)$,
$C(a,D)>0$, and $\rho_c(a,D,\beta)<\frac{1+(K-1)a}{K}$ follows from
$K\bigl(1+(K-1)a-\beta C\bigr)<\bigl(1+(K-1)a\bigr)(K-\beta C)$, which
rearranges to $\beta C\bigl(1-(1+(K-1)a)\bigr)<0$, true since $a<1$.
\end{proof}

This is a statement about where $J_{\rho,\beta}$'s derivative changes sign
along one fixed, hand-specified path (\cref{prop:loss-path}'s affine
$\bm{P}_g(a)$), not a claim about the trajectory of a trained network under
unrestricted SGD, which moves in a much higher-dimensional parameter space
and need not stay near this path at all.
The variable $a$ remains the population-projected coordinate of \cref{prop:loss-path},
held fixed across this analysis as elsewhere. As competence increases, the
process-only reliability requirement $\frac{1+(K-1)a}{K}$ rises. Clean
terminal supervision partially offsets this rise (\cref{eq:rescue-threshold}
is always at most that baseline), but whether it offsets enough to prevent
the net threshold from increasing depends on depth. Near $a=0$,
$\rho_c(a,D,\beta)=\frac1K\bigl[1+(K-1)a\bigr]-\beta\,\frac{D(K-1)^2}{K^2}a^{D-1}+O(a^D)$
for $D\ge3$ (at $D=2$ the correction instead enters at the same order as
the leading term and modifies its coefficient directly): the offset is
higher order in $a$ as $D$ grows, so it matters least at small $a$ for
deep compositions. This local statement does not extend to the whole
interval $(0,1)$: at the paper's own $K=16$ and $\beta=1/D$,
$\rho_c(a,2,\tfrac12)$ is numerically negative for $a\in(0.079,0.421)$,
meaning $J_{\rho,\beta}$ already favors increasing $a$ there even at
$\rho=0$, full process corruption, purely from the clean answer term ---
at this one shallow depth the clean answer alone can carry a wrong rule
toward the true one. For $D\in\{4,6,8\}$, numerical evaluation of
$\rho_c(a,D,1/D)$ shows an increasing threshold over the plotted range.
At $D=3$ it is non-monotone over an intermediate interval: for example,
$\rho_c(0.2,3,1/3)=0.1600$, whereas
$\rho_c(0.3,3,1/3)\approx0.1464$. At each of these depths the threshold
ultimately tends to $1$ as $a\to1$. At $D=8$, for instance, it rises from
$0.06$ at $a=0$ through $0.52$ at $a=0.5$ to $0.90$ at $a=0.9$ and
$0.95$ at $a=0.95$ (\cref{fig:prop8-frontier}). Clean supervision lowers
the threshold relative to the process-only baseline at every $a>0$,
but does not remove the high reliability requirement near $a=1$. This is a barrier to \emph{increasing confidence},
not to \emph{correctness}: along $\bm{P}_g(a)$, the true successor's row
entry is $1/K+a(K-1)/K$ against $(1-a)/K$ for every rival, a margin of
exactly $a$, so greedy decoding from $\bm{P}_g(a)$ already selects the true
successor at every $a>0$ on this path, independent of $\rho$. What
$\rho_c$ governs is instead whether gradient descent along this one path
keeps moving toward higher $a$ at all: at fixed $\rho$, $a$ keeps
increasing locally only while $\rho>\rho_c(a,D,1/D)$. Stationary points
satisfy $\rho_c(a^\star,D,1/D)=\rho$. Non-monotonicity can yield multiple
such points, so a unique stalling point is not guaranteed. Reaching high confidence
at depth $D\ge3$ needs $\rho$ close to $1$, not close to $1/K$, even
though low confidence there is already decodable correctly. One
hypothesis consistent with, but not established by, this diagnostic is
that the measured Boolean-8 frontier's late, steep climb
(\cref{fig:noise-threshold}, between $\rho=0.80$ and $\rho=0.90$ rather
than gradually away from $1/K\approx0.06$) reflects a trained network
needing confident, robust competence rather than merely the bare
existence of a correct greedy decision. Testing this specific hypothesis
would need a direct confidence or margin measurement in the trained
network, which \cref{sec:gradient-alignment}'s diagnostic does not
provide. Three caveats keep this reading honest even as a hypothesis.
First, $\rho_c(a,D,\beta)$ is a stalling threshold along one fixed path,
computed from the objective's derivative at each $a$ separately. It is
not a claim about a trained optimizer's trajectory, and $a$ is not
identified with any of the measured accuracy numbers
above. Second, the sign change is depth-dependent, not monotone the way
this reading suggests: at the one shallow depth $D=2$, the clean answer
cuts the other way and can carry a wrong rule toward the true one from
$\rho=0$, rather than raising the bar. Third, the margin argument above is
specific to this affine path and this population-level table. It says
nothing about the margin realized by a trained Transformer's own
parameters, which need not move along $\bm{P}_g(a)$ at all.
What a trained network's trajectory does on either side of either regime is
not addressed by this static fact.

\begin{figure}[htbp]
 \centering
 \includegraphics[width=0.62\linewidth]{figures/prop8_frontier.pdf}
 \caption{\textbf{The stalling threshold becomes more demanding at high
 competence, with shallow-depth non-monotonicity.} Predicted curves
 (\cref{eq:rescue-threshold}) for $D\in\{2,3,4,6,8\}$ at $K=16$ and
 $\beta=1/D$. Open circles are measured independently from the loss
 definitions (finite differences in $a$, bisection in $\rho$).
 $D\in\{4,6,8\}$ increases over the plotted range. At $D=3$ there is a shallow
 non-monotone region, while $D=2$ dips below zero over an intermediate
 interval. Computed by \texttt{python experiments/run.py prop8-frontier}.
 No model is trained.}
 \label{fig:prop8-frontier}
\end{figure}

This subsection's claims are all about where one fixed path's derivative
changes sign in the shared transition-table model, not about a trained
Transformer.
\Cref{sec:sharp-jump,app:noise-threshold} give the trained-model and
finite-sample numbers separately, as external validation of the
qualitative mechanism rather than verification of this diagnostic's
threshold value.